\section*{Chapter \ref{ch:rc:market-risk-measures} --- Market-Risk Measures}

\begin{solution}{exo:rc:market-risk-measures:1}
$k = \lceil0.01\times500\rceil = 5$: the fifth-largest loss is the VaR. For ES at 97.5\%,
$\lceil0.025\times500\rceil = 13$: the average of the thirteen largest losses.
\end{solution}

\begin{solution}{exo:rc:market-risk-measures:2}
$1.58\text{ million}\times\sqrt{10} = \text{USD}~5.01$ million. The rule fails when changes are
autocorrelated (trends, mean reversion), when volatility changes over the ten days, when the
positions are not held (dynamic hedging) and for non-linear books, whose ten-day P\&L is not ten
one-day P\&Ls.
\end{solution}

\begin{solution}{exo:rc:market-risk-measures:3}
Yellow, with an increase of 0.65: a multiplication factor of 3.65.
\end{solution}

\begin{solution}{exo:rc:market-risk-measures:4}
Alone: $\P(L>0) = 0.9\% < 1\%$, so VaR is 0; the worst 1\% contains the 0.9\% default (100) and 0.1\%
of zeros, ES $= 0.9\times100/1 = 90$. Together: $\P(L>0) = 1-0.991^2 = 1.79\% > 1\%$ while $\P(L>100) =
0.0081\% < 1\%$, so VaR is 100; the worst 1\% holds 0.0081\% of 200 and 0.9919\% of 100, ES $= 100.8$.
\end{solution}

\begin{solution}{exo:rc:market-risk-measures:5}
A short straddle has negative gamma: it loses $\tfrac12\Gamma(\Delta S)^2$ on moves in either
direction. A delta-only revaluation misses that loss, which is largest in exactly the tail
scenarios that set VaR; historical VaR is USD~1.91 million with delta alone and 2.10 with the
gamma term.
\end{solution}

\begin{solution}{exo:rc:market-risk-measures:6}
Tracking volatility sets the scale right on average, but the normal distribution puts too
little probability in the tails: at the 99\% quantile, fat-tailed moves exceed $2.33\hat\sigma$ more
often than 1\% of the time. Exceptions come from the shape, not the level.
\end{solution}

\begin{solution}{exo:rc:market-risk-measures:7}
\texttt{kupiec(1000, 28, 0.01)} gives a likelihood ratio of 22.0 and a p-value of
$3\times10^{-6}$: the model is rejected.
\end{solution}

\begin{solution}{exo:rc:market-risk-measures:8}
Zero exceptions in 250 days happens with probability 8.1\% for a correct model: weak evidence of
conservatism. Over four years the historical model had exactly 1\% exceptions, while the EWMA model,
though lower, failed the \omterm{def:rc:market-risk-measures:tests}{Kupiec test} with 2.8\%. Lower is not better.
\end{solution}

\begin{solution}{pb:rc:market-risk-measures:1}
\textbf{1.} EWMA: VaR USD~1.58 million, ES 1.59 million; historical: 2.10 and 2.34 million.
\textbf{2.} The ten-year yield: 114\% of the EWMA VaR, partly offset by the two-year short.
\textbf{3.} USD~5.01 million.
\textbf{4.} Under normality ES at 97.5\% is 2.338 standard deviations and VaR at 99\% 2.326.
\textbf{5.} The historical figures, or both with an explanation: the board needs the fat-tailed
number and its uncertainty, not the lowest one.
\textbf{6.} 8 exceptions, yellow zone.
\textbf{7.} $3+0.75 = 3.75$.
\textbf{8.} USD~15.02 million at 3, 18.78 million at 3.75, taking today's ten-day VaR for the sixty-day average.
\textbf{9.} None.
\textbf{10.} Yes: with 3 exceptions it would be green with no increase, although its VaR is higher
(USD~2.09 million); the capital comparison depends on the level against the multiplier.
\textbf{11.} Its shape is wrong: the normal tail is too thin for these factors.
\textbf{12.} \omterm{def:rc:market-risk-measures:es}{Expected shortfall} at 97.5\% on a ten-day base horizon scaled by liquidity horizons,
with desk-level backtesting and the P\&L attribution test (chapter~23).
\textbf{13.} 8 exceptions at 99\% is within the 12 allowed; the 97.5\% count would also have to stay
under 30.
\textbf{14.} Whether exceptions cluster: a model can have the right count but fail in bursts,
which is worse for capital adequacy.
\textbf{15.} \omterm{def:rc:market-risk-measures:fhs}{Filtered historical simulation}: fat tails and fast volatility, green over four
years; its VaR is higher, but without a multiplier increase its capital is comparable.
\textbf{16.} A model's job is to be right about the tail; a low number that is exceeded too often
understates capital and misleads limits.
\textbf{17.} Yes: by concentrating risk in positions whose losses sit beyond the quantile (short
options, credit), by trading on stale windows, or by exploiting factors the model omits.
\textbf{18.} ES, stress results, the largest contributions, concentration and liquidity, limit
usage, and backtesting results.
\textbf{19.} Named result: \emph{the quarter-past-four report}: a 99\% VaR of USD~1.58 million and a
97.5\% ES of 1.59 million, 8 exceptions in 250 days, and a multiplier of 3.75.
\textbf{20.} A loss exceeded with a given small probability under recent conditions; not how large
losses are beyond it, nor what happens when conditions change.
\end{solution}

\begin{solution}{iq:rc:market-risk-measures:1}
VaR: the loss quantile at a confidence level over a horizon. ES: the average loss beyond that
quantile. ES sees the tail's size and is coherent (subadditive), so it cannot be improved by
hiding risk beyond the quantile or penalise diversification.
\iqlookfor{the definitions and the tail and coherence arguments.}
\end{solution}

\begin{solution}{iq:rc:market-risk-measures:2}
Historical: real joint moves and fat tails, full revaluation, but slow to react and limited to
the window. Parametric: fast and analytic, linear and normal. Monte Carlo: any distribution and
full revaluation, but model-dependent and costly.
\iqlookfor{trade-offs of assumptions, speed and nonlinearity.}
\end{solution}

\begin{solution}{iq:rc:market-risk-measures:3}
Two independent bonds each defaulting with probability 0.9\%: each has 99\% VaR 0, the pair 100.
\iqlookfor{a concrete counterexample.}
\end{solution}

\begin{solution}{iq:rc:market-risk-measures:4}
Count days whose loss exceeds the previous day's VaR; \omterm{def:rc:market-risk-measures:tests}{Kupiec tests} whether the frequency matches
the level; \omterm{def:rc:market-risk-measures:tests}{Christoffersen tests} whether exceptions are independent over time; the traffic light
maps counts to capital penalties.
\iqlookfor{frequency and independence.}
\end{solution}

\begin{solution}{iq:rc:market-risk-measures:5}
Revalue by trade type with vectorised or grid pricing, sensitivities for linear trades, cache
scenario-independent terms, parallelise by trade and scenario, run incrementally for changed
trades, and keep full revaluation for the non-linear tail.
\iqlookfor{grids, sensitivities and parallelism.}
\end{solution}

\begin{solution}{iq:rc:market-risk-measures:6}
ES is not elicitable alone, but joint VaR--ES tests work: compare realised tail losses with the
predicted ES on exception days, or backtest the VaR at several levels
(the FRTB uses 97.5\% and 99\%).
\iqlookfor{tail-loss comparisons and multi-level VaR tests.}
\end{solution}
